% =====================================================================
%  appendix.tex
% =====================================================================
\appendix
\section{附录}

\subsection{附录 A：文件清单与统计}

\begin{table}[htbp]
\centering
\small
\begin{tabular}{llrr}
\toprule
文件 & 模块 & 原始行数 & 注释后行数\\
\midrule
\file{src/adda.v} & \texttt{adda} & 261 & 417\\
\file{src/rtl.v} & \texttt{rtl} & 80 & 143\\
\file{src/eso.v} & \texttt{eso} & 140 & 252\\
\file{src/divider.v} & \texttt{divider} & 50 & 82\\
\file{src/key.v} & \texttt{key} & 32 & 61\\
\file{src/seg.v} & \texttt{seg} & 153 & 231\\
\file{src/bcd.v} & \texttt{bcd} & 62 & 98\\
\file{src/uart_rx.v} & \texttt{uart\_rx} & 116 & 187\\
\file{src/uart_tx.v} & \texttt{uart\_tx} & 84 & 136\\
\file{src/gowin_pll/gowin_pll.v} & \texttt{Gowin\_PLL} & 137 & 171\\
\file{src/fifo_top/fifo_top.v} & \texttt{fifo\_top} & 504 & 536\\
\file{src/project_2025_11_07.cst} & 约束 & 69 & 109\\
\file{src/gowin_pll/gowin_pll_tmp.v} & 模板 & 19 & 30\\
\file{src/fifo_top/fifo_top_tmp.v} & 模板 & 24 & 34\\
\bottomrule
\end{tabular}
\end{table}

全部注释只增加以 \sig{//} 开头的整行注释，未改动任何代码行。经脚本比对，
注释前后去掉注释后的 Verilog 词法内容\textbf{完全一致}（17 个文件全部通过）。

\subsection{附录 B：模块与关键信号速查}

\begin{table}[htbp]
\centering
\small
\begin{tabular}{lll}
\toprule
模块 & 输入 & 输出\\
\midrule
\texttt{adda} & \sig{ad_date, key_in, rx, sys_clk, sys_rst_n} & \sig{ad_out, sel, seg, tx, led*, \*_clk1}\\
\texttt{rtl} & \sig{ad_data} & \sig{sign, volt}\\
\texttt{eso} & \sig{e_in, u_out, c/k} & \sig{out_da, z1k0..z3k0}\\
\texttt{divider} & \sig{sys_clk} & \sig{clk0, clk1}\\
\texttt{key} & \sig{key_in} & \sig{key_out}\\
\texttt{seg} & \sig{data, point, sign} & \sig{sel, seg}\\
\texttt{bcd} & \sig{data[9:0]} & \sig{unit, ten, hun}\\
\texttt{uart\_rx} & \sig{rx} & \sig{po_data, po_flag}\\
\texttt{uart\_tx} & \sig{pi_data, pi_flag} & \sig{tx}\\
\bottomrule
\end{tabular}
\end{table}

\subsection{附录 C：复现步骤}

\subsubsection{C.1 用安路 TD 打开（代码浏览用容器工程）}

高云 \sig{.gprj} 无法被 TD 直接打开，本报告用 TD 的 Tcl 接口创建了一个同名容器工程：

\begin{lstlisting}[style=cfgcode,numbers=none]
$env:PATH  # 需能找到 D:\FPGA\bin\td_commands_prompt.exe
cd C:/Users/Administrator/Desktop/task/FPGA_TD
# 通过 stdin 喂给 TD 的 Tcl 解释器：
create_project -family ph1 -device PH1A60GEG324 project_2025_11_07_TD .
import_device ph1_60.db -package PH1A60GEG324
add_files -list ./files.list
exit
\end{lstlisting}

随后直接双击 \file{project_2025_11_07_TD.al} 或
\sig{td.exe project_2025_11_07_TD.al} 即可打开，工程只引用 \file{../FPGA/src} 下的源文件，
不复制。注意该工程只能浏览/编辑代码，\textbf{不能}用安路器件综合本设计。

\subsubsection{C.2 编译本报告}

\begin{lstlisting}[style=cfgcode,numbers=none]
cd C:/Users/Administrator/Desktop/task/doc
xelatex main.tex
xelatex main.tex      # 第二遍生成目录与交叉引用
\end{lstlisting}

需要 MiKTeX（含 ctex 宏包与中文字体）。文档引用 \file{../FPGA/src/*.v} 源文件，
移动目录结构后需同步修改相对路径。

\subsection{附录 D：术语表}

\begin{table}[htbp]
\centering
\small
\begin{tabular}{ll}
\toprule
术语 & 含义\\
\midrule
ADRC & 自抗扰控制（Active Disturbance Rejection Control）\\
ESO & 扩张状态观测器（Extended State Observer），估计状态与总扰动\\
TD & 安路科技开发工具 TangDynasty\\
GW5A & 高云 Arora V 系列 FPGA（22nm 工艺）\\
8N1 & UART 8 数据位、无校验、1 停止位\\
FWFT & First-Word Fall-Through，FIFO 首字直通模式\\
BCD & 二进制编码十进制\\
LSB/MSB & 最低有效位 / 最高有效位\\
\bottomrule
\end{tabular}
\end{table}
